\documentclass[11pt]{article}
\usepackage[utf8]{inputenc}
\usepackage{amsmath, amssymb, amsthm}
\usepackage{geometry}
\usepackage{enumitem}
\usepackage{fancyhdr}

\geometry{letterpaper, margin=1in}

\pagestyle{fancy}
\fancyhf{}
\rhead{Math 12.4 - The Cross Product}
\lhead{Practice Problem Set}
\rfoot{Page \thepage}

\title{\textbf{Practice Problem Set: The Cross Product}\\ \large Based on Homework 12.4 Analysis}
\author{}
\date{}

\begin{document}

\maketitle

\section*{Instructions}
Solve the following 30 problems. These questions are designed to test your understanding of the Cross Product, ranging from basic determinants to advanced applications in physics and geometry. 

\textbf{Note:} $\mathbf{i}, \mathbf{j}, \mathbf{k}$ represent the standard basis vectors. $\langle x, y, z \rangle$ represents vector components.

\hrulefill

\section*{Part 1: Computational Fluency}

\begin{enumerate}
    \item \textbf{Standard Calculation:} Given $\mathbf{a} = \langle 3, -2, 1 \rangle$ and $\mathbf{b} = \langle 1, 4, -2 \rangle$, calculate $\mathbf{a} \times \mathbf{b}$.
    
    \item \textbf{Basis Vectors:} Calculate $(3\mathbf{i} - \mathbf{j} + 2\mathbf{k}) \times (\mathbf{i} + 5\mathbf{k})$.
    
    \item \textbf{Variable Components:} Let $\mathbf{u} = \langle t, 1, 0 \rangle$ and $\mathbf{v} = \langle t^2, -t, 2 \rangle$. Find $\mathbf{u} \times \mathbf{v}$ in terms of $t$.
    
    \item \textbf{Trigonometric Identity:} Let $\mathbf{a} = \langle \sin\theta, \cos\theta, 0 \rangle$ and $\mathbf{b} = \langle -\cos\theta, \sin\theta, 1 \rangle$. Compute $\mathbf{a} \times \mathbf{b}$ and simplify using Pythagorean identities.
    
    \item \textbf{Parallel Check:} Determine if $\mathbf{a} = \langle -2, 3, 5 \rangle$ and $\mathbf{b} = \langle 4, -6, -10 \rangle$ are parallel, perpendicular, or neither. Justify using the cross product.

\section*{Part 2: Algebraic Properties}
    
    \item \textbf{Simplification:} Simplify the expression $2\mathbf{i} \times (\mathbf{j} - 3\mathbf{k})$ using the cyclic properties of basis vectors.
    
    \item \textbf{Distributivity:} Expand and simplify $(\mathbf{a} + 2\mathbf{b}) \times (2\mathbf{a} - \mathbf{b})$. Express your answer in terms of $(\mathbf{a} \times \mathbf{b})$.
    
    \item \textbf{Order of Operations:} Evaluate $(\mathbf{j} \times \mathbf{i}) \times \mathbf{k}$. Caution: Remember vector triple product order matters.
    
    \item \textbf{Orthogonality Verification:} If $\mathbf{c} = \mathbf{a} \times \mathbf{b}$, show that $\mathbf{c} \cdot \mathbf{a} = 0$ for vectors $\mathbf{a} = \langle 1, 5, 2 \rangle$ and $\mathbf{b} = \langle -2, 1, 3 \rangle$.
    
    \item \textbf{Anticommutativity:} True or False: $(\mathbf{a} \times \mathbf{b}) \cdot \mathbf{c} = - (\mathbf{b} \times \mathbf{a}) \cdot \mathbf{c}$. Explain.

\section*{Part 3: Geometric Applications (Area and Angle)}
    
    \item \textbf{Parallelogram Area:} Find the area of the parallelogram spanned by vectors $\mathbf{u} = \langle 2, 0, 1 \rangle$ and $\mathbf{v} = \langle 0, 4, 3 \rangle$.
    
    \item \textbf{Triangle Area (3D):} Find the area of the triangle with vertices $P(1, 0, 1)$, $Q(2, 2, 2)$, and $R(3, -1, 4)$.
    
    \item \textbf{2D Polygon in 3D:} Points $A(0,0)$, $B(4,1)$, and $C(2,5)$ lie in the $xy$-plane. Use the cross product (treating $z=0$) to find the area of triangle $ABC$.
    
    \item \textbf{Angle Extraction:} Given $|\mathbf{a} \cdot \mathbf{b}| = 10\sqrt{3}$ and $|\mathbf{a} \times \mathbf{b}| = 10$, find the angle $\theta$ between vectors $\mathbf{a}$ and $\mathbf{b}$ where $0 \le \theta \le \pi$.
    
    \item \textbf{Lagrange's Identity Application:} If $|\mathbf{a}| = 5$, $|\mathbf{b}| = 4$, and $\mathbf{a} \cdot \mathbf{b} = 12$, find $|\mathbf{a} \times \mathbf{b}|$ without finding the vectors.

\section*{Part 4: Triple Products (Scalar and Vector)}
    
    \item \textbf{Volume of Parallelepiped:} Find the volume of the parallelepiped determined by vectors $\mathbf{a} = \langle 1, 2, 3 \rangle$, $\mathbf{b} = \langle -1, 1, 2 \rangle$, and $\mathbf{c} = \langle 2, 1, 4 \rangle$.
    
    \item \textbf{Coplanar Vectors:} Find the value of $x$ such that the vectors $\langle 2, 1, 0 \rangle$, $\langle 1, -1, 2 \rangle$, and $\langle 0, x, 4 \rangle$ are coplanar.
    
    \item \textbf{Vector Triple Product (BAC-CAB):} Compute $\mathbf{i} \times (\mathbf{j} \times \mathbf{i})$ using the formula $\mathbf{a} \times (\mathbf{b} \times \mathbf{c}) = \mathbf{b}(\mathbf{a} \cdot \mathbf{c}) - \mathbf{c}(\mathbf{a} \cdot \mathbf{b})$.
    
    \item \textbf{Cyclic Scalar Triple Product:} True or False: $\mathbf{a} \cdot (\mathbf{b} \times \mathbf{c}) = \mathbf{b} \cdot (\mathbf{c} \times \mathbf{a})$. 
    
    \item \textbf{Mixed Product Computation:} Given $\mathbf{a} \perp \mathbf{b}$, $|\mathbf{a}|=3, |\mathbf{b}|=2$. What is $|\mathbf{a} \times (\mathbf{a} \times \mathbf{b})|$?

\section*{Part 5: Applications and Modeling}
    
    \item \textbf{Torque Calculation:} A bolt is tightened by applying a 40 N force to a 0.25 m wrench. The force vector is $\mathbf{F} = \langle 0, 30, -10 \rangle$ and the position vector of the handle is $\mathbf{r} = \langle 0.25, 0, 0 \rangle$. Find the magnitude of the torque $\tau = \mathbf{r} \times \mathbf{F}$.
    
    \item \textbf{Distance Point to Line:} Use the formula $d = \frac{|\vec{QP} \times \mathbf{u}|}{|\mathbf{u}|}$ to find the distance from point $P(2,1,3)$ to the line passing through $Q(0,0,0)$ with direction $\mathbf{u} = \langle 1, 2, -2 \rangle$.
    
    \item \textbf{Unit Normal Vector:} Find a unit vector orthogonal to both $\mathbf{a} = \langle 1, -1, 0 \rangle$ and $\mathbf{b} = \langle 1, 1, 1 \rangle$ pointing in the positive $z$ direction.
    
    \item \textbf{Financial Volatility Surface:} A financial model defines a surface position $\mathbf{P}(k,t)$. Tangent vectors are found to be $\mathbf{T}_k = \langle 1, 0, 0.5 \rangle$ and $\mathbf{T}_t = \langle 0, 1, -0.2 \rangle$. Calculate the normal vector $\mathbf{N} = \mathbf{T}_k \times \mathbf{T}_t$.
    
    \item \textbf{Magnetic Force:} The force on a moving charge is $\mathbf{F} = q(\mathbf{v} \times \mathbf{B})$. If a proton ($q > 0$) moves in the $+x$ direction ($\mathbf{v}$) through a magnetic field pointing in the $+y$ direction ($\mathbf{B}$), in which direction is the force?

\section*{Part 6: Conceptual Diagnostics ("Find the Flaw")}
    
    \item \textbf{Valid Operations:} Which of the following expressions is undefined? 
    (a) $(\mathbf{a} \times \mathbf{b}) \times \mathbf{c}$ 
    (b) $(\mathbf{a} \cdot \mathbf{b}) \times \mathbf{c}$ 
    (c) $|\mathbf{a}|(\mathbf{b} \times \mathbf{c})$.
    
    \item \textbf{Determinant Errors:} A student calculates $\mathbf{a} \times \mathbf{b}$ for $\mathbf{a}=\langle 1, 2, 0 \rangle, \mathbf{b}=\langle 3, 4, 0 \rangle$ and gets $\mathbf{i}(0) + \mathbf{j}(0) + \mathbf{k}(4-6) = -2\mathbf{k}$. Did they make a "middle minus" sign error?
    
    \item \textbf{Self-Cross:} Evaluate $(\mathbf{a} + \mathbf{b}) \times (\mathbf{a} + \mathbf{b})$.
    
    \item \textbf{Right Hand Rule:} $\mathbf{u}$ points North. $\mathbf{v}$ points West. In what geographic direction does $\mathbf{u} \times \mathbf{v}$ point? (Up or Down).
    
    \item \textbf{Triangle Inequality:} Is it true that $|\mathbf{a} \times \mathbf{b}| \le |\mathbf{a}| |\mathbf{b}|$? Justify with the geometric definition.

\end{enumerate}

\newpage

\section*{Answer Key and Solutions}

\begin{enumerate}
    \item $\langle 0, 7, 14 \rangle$.
    \item $\langle -5, -13, 1 \rangle$. (Hint: Distribute or use determinant).
    \item $\langle 2, -2t, -2t^2 \rangle$.
    \item $\langle \cos\theta, -\sin\theta, 1 \rangle$. Note: $\sin^2\theta + \cos^2\theta = 1$ simplifies the $z$-component.
    \item \textbf{Parallel.} $\mathbf{b} = -2\mathbf{a}$, so $\mathbf{a} \times \mathbf{b} = \mathbf{0}$.
    \item $2\mathbf{k} + 6\mathbf{j}$. Recall $\mathbf{i} \times \mathbf{j} = \mathbf{k}$ and $\mathbf{i} \times \mathbf{k} = -\mathbf{j}$.
    \item $-5(\mathbf{a} \times \mathbf{b})$. (Expansion: $2\mathbf{a}\times\mathbf{a} - \mathbf{a}\times\mathbf{b} + 4\mathbf{b}\times\mathbf{a} - 2\mathbf{b}\times\mathbf{b} = 0 - \mathbf{a}\times\mathbf{b} - 4\mathbf{a}\times\mathbf{b} - 0$).
    \item $\mathbf{0}$. $\mathbf{j} \times \mathbf{i} = -\mathbf{k}$. Then $-\mathbf{k} \times \mathbf{k} = \mathbf{0}$.
    \item Dot product $\langle 13, -7, 11 \rangle \cdot \langle 1, 5, 2 \rangle = 13 - 35 + 22 = 0$.
    \item \textbf{True.} $\mathbf{b} \times \mathbf{a} = -(\mathbf{a} \times \mathbf{b})$. The negative sign pulls out.
    \item Magnitude is $\sqrt{(-4)^2 + (-6)^2 + 8^2} = \sqrt{16+36+64} = \sqrt{116} = 2\sqrt{29}$.
    \item $\vec{PQ} = \langle 1, 2, 1 \rangle, \vec{PR} = \langle 2, -1, 3 \rangle$. Cross = $\langle 7, -1, -5 \rangle$. Area = $\frac{\sqrt{75}}{2} = \frac{5\sqrt{3}}{2}$.
    \item $\vec{AB}=\langle 4, 1, 0 \rangle, \vec{AC}=\langle 2, 5, 0 \rangle$. Cross $\langle 0, 0, 18 \rangle$. Area = 9.
    \item $\tan\theta = \frac{|\mathbf{a}\times\mathbf{b}|}{\mathbf{a}\cdot\mathbf{b}} = \frac{10}{10\sqrt{3}} = \frac{1}{\sqrt{3}}$. $\theta = 30^\circ$ or $\pi/6$.
    \item $|\mathbf{a}\times\mathbf{b}|^2 = |\mathbf{a}|^2|\mathbf{b}|^2 - (\mathbf{a}\cdot\mathbf{b})^2 = (25)(16) - 144 = 400 - 144 = 256$. Magnitude $= 16$.
    \item Scalar triple product determinant = $-5$. Volume = $|-5| = 5$.
    \item Determinant of matrix rows is $-8 + 3x + 8 = 3x$. For coplanar, Volume $= 0 \implies 3x=0 \implies x=0$.
    \item $\mathbf{i} \times (-\mathbf{k}) = \mathbf{j}$. Formula: $\mathbf{j}(\mathbf{i}\cdot\mathbf{i}) - \mathbf{i}(\mathbf{i}\cdot\mathbf{j}) = \mathbf{j}(1) - 0 = \mathbf{j}$.
    \item \textbf{True.} Cyclic permutation of scalar triple product preserves value.
    \item Since $\mathbf{a} \perp \mathbf{b}$, $|\mathbf{a} \times \mathbf{b}| = 6$. The vector $\mathbf{v} = \mathbf{a} \times \mathbf{b}$ is perp to $\mathbf{a}$. So $|\mathbf{a} \times \mathbf{v}| = |\mathbf{a}||\mathbf{v}|\sin(90^\circ) = 3 \cdot 6 \cdot 1 = 18$.
    \item $\tau = \langle 0, 2.5, 7.5 \rangle$. Magnitude $\approx 7.91$ Nm.
    \item $\vec{QP} = \langle 2, 1, 3 \rangle$. Cross product with direction: $\langle -8, 5, 3 \rangle$. Mag $\sqrt{98}$. Direction mag $|\mathbf{u}| = 3$. Distance $\frac{\sqrt{98}}{3} = \frac{7\sqrt{2}}{3}$.
    \item $\mathbf{a} \times \mathbf{b} = \langle -1, -1, 2 \rangle$. Magnitude $\sqrt{6}$. Unit vector $\langle \frac{-1}{\sqrt{6}}, \frac{-1}{\sqrt{6}}, \frac{2}{\sqrt{6}} \rangle$.
    \item $\langle -0.5, 0.2, 1 \rangle$. (Direct determinant calculation).
    \item $\mathbf{v}=\langle v,0,0 \rangle, \mathbf{B}=\langle 0, B, 0 \rangle$. $\mathbf{v} \times \mathbf{B} = vB(\mathbf{k})$. Direction is $+\mathbf{z}$ (Out of page).
    \item \textbf{(b)} is undefined. Scalar crossed with vector.
    \item No, the student is correct. The middle term is $-\mathbf{j}(0-0) = 0$. The error usually happens when the middle determinant is non-zero.
    \item $\mathbf{0}$. Any vector crossed with itself is the zero vector.
    \item $\mathbf{u}=\mathbf{j}$ (North implies $y$), $\mathbf{v}=-\mathbf{i}$ (West implies $-x$). $\mathbf{j} \times (-\mathbf{i}) = -(\mathbf{j} \times \mathbf{i}) = -(-\mathbf{k}) = \mathbf{k}$. Direction: \textbf{Up} (away from ground).
    \item \textbf{True.} $|\mathbf{a} \times \mathbf{b}| = |\mathbf{a}||\mathbf{b}|\sin\theta$. Since $|\sin\theta| \le 1$, the inequality holds.
\end{enumerate}

\newpage

\section*{Concept Checklist \& Map}

This checklist maps the specific concepts and problem types from the source text to the questions above. Use this to verify your coverage.

\begin{table}[h]
\centering
\begin{tabular}{|l|l|l|}
\hline
\textbf{Concept / Problem Type} & \textbf{Source Section} & \textbf{Problem \#} \\ \hline
Determinant Calculation (Integers) & Part 1.4, Part 2 (Prob 1) & 1, 5, 9 \\ \hline
Basis Vectors \& Zero Comps & Part 2 (Prob 3) & 2, 6 \\ \hline
Variable/Function Components & Part 2 (Prob 6, 7) & 3, 4 \\ \hline
Algebraic Distributivity & Part 2 (Prob 11), Part 3 & 7 \\ \hline
Cyclic Basis Properties & Part 2 (Prob 11), Part 7 & 6, 8, 18 \\ \hline
Orthogonality Check (Dot=0) & Part 1.2, Part 5 & 9 \\ \hline
Anticommutativity & Part 3 (Pitfalls), Part 8 & 10 \\ \hline
Parallelogram Area & Part 1.4, Part 2 (Prob 13) & 11 \\ \hline
Triangle Area (3D points) & Part 2 (Prob 16, 28) & 12 \\ \hline
2D Area (using z=0) & Part 4.2 (Tricks) & 13 \\ \hline
Angle from Mag/Dot Ratio & Part 2 (Prob 19) & 14 \\ \hline
Lagrange Identity ($|\mathbf{a}\times\mathbf{b}|^2$) & Part 1.2 (Derivation) & 15 \\ \hline
Scalar Triple Product (Volume) & Part 1.4, Part 2 (Prob 33) & 16, 19 \\ \hline
Coplanar Check (Vol=0) & Part 2 (Prob 38) & 17 \\ \hline
Vector Triple Product & Part 7.1 (Missing Concept) & 18, 20 \\ \hline
Physics: Torque & Part 6.1 & 21 \\ \hline
Distance Point to Line & Part 7.2 (Missing Concept) & 22 \\ \hline
Normal Vectors (Planes/Surfaces) & Part 6.2 & 23, 24 \\ \hline
Physics: Magnetic Force (RHR) & Part 1.3 & 25 \\ \hline
"Find the Flaw" / Validity & Part 8 & 26, 27, 28 \\ \hline
Right Hand Rule & Part 1.1, Part 2 (Prob 14) & 29 \\ \hline
Inequalities/Definitions & Part 1.1 & 30 \\ \hline
\end{tabular}
\caption{Cross reference of problems to source PDF concepts.}
\end{table}

\end{document}